\section{讨论}

\subsection{三种方法的精度与效率}

若以应力精度为主要目标，Q4 是三种方法中最均衡的选择。它在 $160\times32$ 网格上的探针应力误差不超过 $0.042\%$，明显低于 CST 的 $0.146\%$，且能直接给出单元内线性变化的应力场；代价是每节点两个自由度、单元数虽少但后处理更复杂。CST 的优势在实现简单与几何适应性，常应变假设使其偏刚，局部应力误差更大，且在粗网格上呈现明显的分块与对角线敏感性；若必须使用 CST，应配合面积加权节点恢复与更密网格。

杆单元的效率无可替代：每节点一个自由度，节点位移精确，且与网格无关。它的局限是只给出一维曲线，无法描述端部局部变形、横向应力与截面非均匀性；因此它适合 A/B 级工况的快速估算与趋势判断，不适合作为 C 级工况或端部强度评定的最终依据。工程上合理的流程是“一维杆快速估算 + Q4 二维复核 + 关键区域网格加密”，CST 则用于复杂几何的快速原型或与自动剖分工具对接。

\subsection{一维判据与边界层}

$\eta=|\beta|h_0/L=\tan\alpha$ 是一个纯几何判据，优点是可在建模前计算，缺点是没有包含端部约束与载荷分布的影响。本文用二维参考标定了 A/B/C 阈值，并给出对应的后验偏差；需要强调两点。第一，即使 $\eta=0$（等截面），固定端全约束仍造成约 $0.25\%$–$0.50\%$ 的位移偏差，这是圣维南边界层而非锥度效应，一维杆的截面平均应力仍精确等于 $P/A$。第二，A 级只承诺“内部/中截面可信、端部约 $1\%$”，不应把 A 级理解为“端部位移也准确”；若关心端部，必须做二维复核。

\subsection{口径敏感性}

本文量化了三类口径的影响，结论是：端载口径影响端部位移约 $0.02\%$–$0.28\%$，但影响加载端局部应力可达 $1\%$–$8\%$；截面平均口径（节点算术平均与面积加权）在粗网格上可差 $0.4\%$，并严重污染位移收敛阶；CST 后处理的 raw 与 recovered 差异主要出现在粗网格的截面内振荡。因此，跨方法/跨代码对比前必须统一端载、截面平均和应力恢复口径，否则很容易把后处理差异误判为方法误差。

\subsection{局限}

本文限定于线弹性、小变形、平面应力、对称楔形与端部均布载荷；未考虑三维效应、材料非线性、接触、动力与优化。CalculiX 平面应力单元通过三维展开实现，与严格二维平面应力在厚度积分细节上存在约 $0.01\%$–$0.05\%$ 的差异。CST 的对角线敏感性指标依赖镜像对称性；对非对称几何/载荷，全局量也会出现离散，需要按一般网格敏感性方法处理。此外，本文的“梯度保真误差”使用截面平均应力的中心差分，参考梯度来自 CPS8 细网格，其绝对值仍受参考解与差分算子影响，宜作为相对比较而非绝对误差。
